(1) 旋转系假想加速度为 \(-2\boldsymbol\Omega\times\mathbf v\)。直接计算叉乘得 \(a_x=-2\Omega(\cos\lambda v_z-\sin\lambda v_y)\)、\(a_y=-2\Omega\sin\lambda v_x\)、\(a_z=2\Omega\cos\lambda v_x\)。此处 \(g\) 已含零阶离心效应，忽略 \(O(\Omega^2)\)。
(2) 代入零阶速度 \((v_0,0,w_0-gt)\)，并以初始位置为零积分两次，得到 \(x=v_0t-\Omega\cos\lambda(w_0t^2-gt^3/3)\)、\(y=-\Omega v_0t^2\sin\lambda\)、\(z=w_0t-gt^2/2+\Omega v_0t^2\cos\lambda\)，式中误差均为 \(O(\Omega^2)\)。
(3) 将 \(z(t_f)=0\) 的非零根展开，\(t_f=2w_0/(g-2\Omega v_0\cos\lambda)=2w_0/g+4\Omega v_0w_0\cos\lambda/g^2+O(\Omega^2)\)。求 \(y\) 的一阶量时只需代入零阶时间，得到 \(y_f=-4\Omega v_0w_0^2\sin\lambda/g^2\)。
(4) 赤道处 \(\sin\lambda=0\)，一阶南北偏移为零，但竖向加速度可变化；北极 \(\cos\lambda=0\) 时竖向修正消失而南北偏移最大。必须有 \(\Omega t_f\ll1\)，且假设曲率、风阻及重力场随高度的变化可忽略。
